\Needspace{10\baselineskip}
\begin{center}
PALLAS.
\end{center}


Monatliche Correspondenz. Band VII, S.\ 369—375 und 466—469. April und Mai 1803.


[Das Aprilheft gibt den Inhalt des Briefs an \textsc{Bode} vom 3.\ März wieder. Siehe oben S.\ 233. Das Maiheft enthält noch folgende Stelle:]

— — — Diese Beobachtungen [von Dr.\ \textsc{Olbers}] stimmen weiterhin ausgezeichnet mit den Elementen VI der Bahn von Dr.\ \textsc{Gauß} überein. Hier ist der Vergleich:

\begin{center}
\begingroup\fontsize{10}{12}\selectfont\setlength{\tabcolsep}{3pt}\renewcommand{\arraystretch}{1.12}
\begin{tabular}{lrl}
\multicolumn{3}{l}{1803 März 4. $17^u\,11^m\,41^s$ m. Z. Br.} \\
Berechn. AR. $\PallasSymbol$ & $275^\circ\,52'\,35''9$ & Unterschied. R—B, $-2''1$ \\
Berechn. Decl. $\PallasSymbol$ & $8\ 58\ 18.1$ & Unterschied. R—B, $-4.9$ \\
\end{tabular}\endgroup
\end{center}

Da der Fehler der vorhergehenden Elemente V so klein und außerdem immer hinreichend genau bekannt ist, wird die in unserem Dezemberheft angegebene Ephemeride während ihres gesamten Zeitraums zum Auffinden des Planeten vollkommen ausreichen. Die Rektaszension darin ist jetzt $2\tfrac{1}{2}$ Min.\ zu groß und wird am 9.\ August $4'\,6''$ zu groß sein, sofern man sich auf die fortdauernde Übereinstimmung der neuen Element verlassen kann. — — — Dr.\ \textsc{Gauß} hat inzwischen die Ephemeride der \emph{Pallas} [in der von ihm bisher verwendeten Form] nach den neuen Elementen VI [vom 9.\ Aug.] bis zum 23.\ Okt.\ 1803 fortgesetzt. Länger wird man sie wohl nicht verfolgen können, da sie schon am 18.\ Sept.\ lichtschwächer zu werden beginnt, als sie bei ihrer Wiederentdeckung war.

\begin{flushright}
[\textsc{von Zach}.]
\end{flushright}

